\documentclass[11pt,a4paper]{article}
\usepackage[T1]{fontenc}
\usepackage[utf8]{inputenc}
\usepackage{lmodern,microtype}
\usepackage[margin=25mm]{geometry}
\usepackage{amsmath,amssymb,amsthm,booktabs,enumitem}
\usepackage{needspace}
\usepackage[hidelinks]{hyperref}
\hypersetup{pdftitle={Exact Natarajan Dimension of the Natural PSL(2,q) Action},pdfauthor={Ryutaro Yonezu}}
\usepackage{fancyhdr}
\setlength{\headheight}{14pt}
\pagestyle{fancy}
\fancyhf{}
\lhead{\small Natarajan dimension of the natural $\mathrm{PSL}(2,q)$ action}
\rhead{\small No.42 / v1.0}
\cfoot{\thepage}
\newtheorem{theorem}{Theorem}
\newtheorem{lemma}{Lemma}
\newtheorem{proposition}{Proposition}
\DeclareMathOperator{\PSL}{PSL}
\DeclareMathOperator{\PGL}{PGL}
\DeclareMathOperator{\Sym}{Sym}
\DeclareMathOperator{\Ndim}{Ndim}
\newcommand{\F}{\mathbb{F}}
\newcommand{\PP}{\mathbb{P}}
\setlist{nosep}
\title{\vspace{-12mm}\Large\bfseries Exact Natarajan Dimension of the\protect\\ Natural $\mathrm{PSL}(2,q)$ Action}
\author{Ryutaro Yonezu\\\small Yonezu Research Lab\\\small\texttt{yonezaemon1+research@gmail.com}}
\date{\small Research No.42 --- Version 1.0 --- 27 October 2026}
\begin{document}
\maketitle
\thispagestyle{plain}
\begin{abstract}
Let $\mathcal H_q$ be the class of permutations of $\PP^1(\F_q)$ induced by
the natural action of $\PSL(2,q)$, where $q$ ranges over prime powers. We give
an elementary proof that its Natarajan dimension is $1$ for $q=2$, is $2$
for $q\in\{3,4,5,7,11\}$, and is $3$ for every other prime power.
For odd $q$, a normalized three-dimensional Natarajan cube is characterized
by quadratic-character conditions. When $q\equiv1\pmod4$, a four-cycle in
a nonneighbor-induced Paley graph supplies a cube; when $q\equiv3\pmod4$,
cube existence is equivalent to a transitive five-vertex Paley subtournament.
The proof covers all finite fields, with finite computations supplied as
checks. The ingredients and several small cases are known. Novelty of the
exact all-$q$ classification has not been established; no priority claim is made.
\end{abstract}
\noindent\textbf{Status.} Unreviewed preprint. AI-assisted preparation and rechecking do not
constitute independent third-party review or human peer review.

\section{Model, scope, and result}
Put $\Omega_q=\PP^1(\F_q)=\F_q\cup\{\infty\}$ and
$\mathcal H_q=\{x\mapsto gx:g\in\PSL(2,q)\}$ on its full domain.
An invertible matrix $\begin{pmatrix}a&b\\c&d\end{pmatrix}$ acts as
$x\mapsto (ax+b)/(cx+d)$, with the usual conventions at poles and at infinity.
It represents an element of $\PSL(2,q)$ precisely when its determinant is a
nonzero square. This follows by rescaling it to determinant one. In even
characteristic every nonzero element is a square, so $\PSL(2,q)=\PGL(2,q)$.

A set of distinct points $x_1,\ldots,x_m$ is \emph{Natarajan-shattered} if
there are labels $u_i\ne v_i$ such that every choice from
$\prod_{i=1}^m\{u_i,v_i\}$ is realized by a member of $\mathcal H_q$.
Let $N(q)=\Ndim(\mathcal H_q)$ be the maximum possible $m$.
This is the standard multiclass dimension, not a new definition.

\begin{theorem}\label{thm:main}
For every prime power $q$,
\[
N(q)=\begin{cases}
1,&q=2,\\
2,&q\in\{3,4,5,7,11\},\\
3,&\text{every other prime power }q.
\end{cases}
\]
\end{theorem}

The statement concerns this specific action of degree $q+1$, not all actions
of the abstract group. For example, $\PSL(2,4)$ and $\PSL(2,5)$ are both
isomorphic to $A_5$, but act here on five and six points, respectively.
The result does not assert improved ARC scores or certification of an
unrestricted ARC task's intended rule.

\Needspace{10\baselineskip}
\section{Elementary reductions}
\begin{lemma}\label{lem:labels}
For a class of injective maps, the output pairs of a Natarajan-shattered
set are pairwise disjoint. Thus shattering $m$ inputs requires at least $2m$ labels.
\end{lemma}
\begin{proof}
If two different inputs had a common available label, the pattern selecting
that label at both would not be injective. Within a single pair the labels
are different by definition.
\end{proof}

\begin{lemma}\label{lem:bounds}
For all $q$, $N(q)\le3$. For $q\ge3$, $N(q)\ge2$.
\end{lemma}
\begin{proof}
Two projective linear transformations agreeing at three distinct points
are identical. A shattered four-set would give two realizers agreeing at
three inputs and disagreeing at the fourth, a contradiction.

The natural $\PSL(2,q)$ action is doubly transitive. Indeed, translations
and $x\mapsto-1/x$ are represented by determinant-one matrices and make
the action transitive; the stabilizer of infinity contains all translations,
which are transitive on the remaining points. For $q\ge3$ there are at least
four labels. Choose two disjoint output pairs; double transitivity realizes
all four choices on two fixed inputs.
\end{proof}

For $q=2$, Lemma~\ref{lem:labels} excludes dimension two, and a moved point
gives dimension one. For $q=3,4$, there are fewer than six labels, so the two
lemmas give dimension two. In particular, $q=3$ already has four projective
points and is covered by the lower-bound argument without an extra exception.

\section{A normalized cube criterion for odd \texorpdfstring{$q$}{q}}\label{sec:criterion}
Assume $q$ is odd. Any shattered triple can be normalized so that the inputs
are $(\infty,0,1)$ and one output at each input equals that input.
First compose all realizers on the left by the inverse of the all-zero
realizer. Then conjugate the whole action by the unique $\PGL(2,q)$ map
sending the three inputs to $(\infty,0,1)$. Such conjugation preserves
$\PSL(2,q)$: it is the kernel of the determinant-square-class homomorphism
on $\PGL(2,q)$.

Lemma~\ref{lem:labels} now makes the output pairs
\[
\{\infty,a\},\qquad\{0,b\},\qquad\{1,c\},
\]
where $a,b,c,0,1$ are five distinct finite elements. Let $\chi$ denote the
quadratic character of $\F_q$, extended by $\chi(0)=0$.

\begin{proposition}\label{prop:criterion}
These pairs form a Natarajan cube exactly when
\begin{align}
\chi(c)=\chi(1-b)=\chi(c-b)&=1,\label{eq:cross}\\
\chi(a)=\chi(a-1)=\chi(a-b)&=\chi(a-c).\label{eq:uniform}
\end{align}
\end{proposition}
\begin{proof}
The projective map sending $(\infty,0,1)$ to $(\infty,y,z)$ is
$x\mapsto(z-y)x+y$. Its determinant is $z-y$, so membership in $\PSL(2,q)$
is equivalent to $\chi(z-y)=1$. For $y\in\{0,b\}$ and $z\in\{1,c\}$,
the four conditions are exactly \eqref{eq:cross}, together with $\chi(1)=1$.

The unique projective map sending $(\infty,0,1)$ to $(a,y,z)$ is represented by
\[
\begin{pmatrix}
a(z-y)&y(a-z)\\
z-y&a-z
\end{pmatrix}.
\]
Its determinant is $(a-y)(a-z)(z-y)$. All factors are nonzero. Under
\eqref{eq:cross}, its square class is positive exactly when
$\chi(a-y)=\chi(a-z)$. Taking all four choices of $y,z$ gives
\eqref{eq:uniform}. Thus all eight patterns are accounted for in both directions.
\end{proof}

The triple-character invariant behind this calculation is known; see
Kr\v{c}adinac, Pav\v{c}evi\'{c}, and Tabak~\cite[Lemma~2.2]{kpt}.
We do not claim that invariant or the determinant-square test as new.

\section{Character sums and the case \texorpdfstring{$q\equiv1\pmod4$}{q = 1 (mod 4)}}
We use the standard identities
\begin{equation}\label{eq:sums}
\sum_{x\in\F_q}\chi(x)=0,\qquad
\sum_{x\in\F_q}\chi(x)\chi(x-t)=-1\quad(t\ne0).
\end{equation}
For completeness, scale to $t=1$ in the second sum. The number of pairs
$(x,y)$ with $y^2=x(x-1)$ is $q$ plus that sum. Put $U=2x-1$, $V=2y$;
the equation becomes $(U-V)(U+V)=1$. The choices $(U-V,U+V)=(r,r^{-1})$
for $r\ne0$ give exactly $q-1$ solutions, proving the identity. The first
sum follows by multiplication by a nonsquare. Also $\chi(-1)=1$ precisely
when $q\equiv1\pmod4$.

Make the Paley graph on $\F_q$ by joining $x,y$ when $\chi(y-x)=1$.
For $q\equiv1\pmod4$ this is an undirected graph. Let
$S=\{x:\chi(x)=-1\}$, the nonneighbors of zero. Its size is
$m=(q-1)/2$. The induced graph on $S$ is $k$-regular with $k=(q-1)/4$.
To check the degree directly, fix $x\in S$ and expand
\[
\frac14\sum_{y\in\F_q}(1-\chi(y))(1+\chi(y-x)).
\]
Equation~\eqref{eq:sums} gives $(q+1)/4$. The $y=0$ term is zero, while
the $y=x$ term is $1/2$ and must be removed. The actual degree is $(q-1)/4$.

If $q\ge13$, then $k\ge3$ and $m=2k$. A $C_4$-free graph has at most one
common neighbor for each pair of vertices. Counting two-edge paths gives
\[
m\binom{k}{2}\le\binom{m}{2},\qquad k(k-1)\le2k-1,
\]
which is impossible for integer $k\ge3$. Thus the induced graph contains
a four-cycle, not necessarily an induced cycle. For $q=9$, it is a
two-regular simple graph on four vertices, so it too is a four-cycle.

Write the cycle's bipartition as $\{u,u'\}$ and $\{v,v'\}$. All four cross
edges are present. The affine map $h(t)=(t-u)/(v-u)$ has a square multiplier.
Set $a=h(0)$, $b=h(u')$, $c=h(v')$. The cross edges give
\eqref{eq:cross}, and the common nonadjacency to zero gives
\eqref{eq:uniform} with common value $-1$. The five finite elements are
distinct. Proposition~\ref{prop:criterion} gives $N(q)=3$ for all
$q\equiv1\pmod4$ with $q\ge9$.

For $q=5$, any $a$ has only two other elements with $\chi(a-x)=1$ and two
with value $-1$. Four distinct elements cannot all have the same value as
required by \eqref{eq:uniform}. Hence $N(5)=2$.

\section{The case \texorpdfstring{$q\equiv3\pmod4$}{q = 3 (mod 4)}}
Make the Paley tournament with arc $x\to y$ when $\chi(y-x)=1$.
Equations~\eqref{eq:cross}--\eqref{eq:uniform} are equivalent to existence
of a transitive five-vertex subtournament.

For the forward implication, all four cross arcs run from $\{0,b\}$ to
$\{1,c\}$. Each pair has an internal order, making the four vertices
transitive. Equation~\eqref{eq:uniform} makes $a$ a common source or sink,
giving a transitive five-set.

Conversely, order a transitive five-set as $t_1\to t_2\to t_3\to t_4\to t_5$,
with every arc pointing forward. Take $a_0=t_5$, $\{u,u'\}=\{t_1,t_2\}$,
and $\{v,v'\}=\{t_3,t_4\}$. The map $h(t)=(t-u)/(v-u)$ has square
multiplier. Then $a=h(a_0)$, $b=h(u')$, $c=h(v')$ satisfy
\eqref{eq:cross} and \eqref{eq:uniform}; the last vertex is a common sink.

Every pair of Paley tournament vertices has exactly $(q-3)/4$ common
outneighbors. To see this, expand the two positive-character indicators
and use \eqref{eq:sums}. The sum before deleting the two endpoints is
$(q-1)/4$; the endpoint contributions add to $1/2$, since $\chi(-1)=-1$.
The result is $(q-3)/4$, independently of the arc direction.

For $q=7,11$, that number is respectively one and two. The first two
vertices of a transitive five-set would have three common outneighbors,
which is impossible. Thus $N(7)=N(11)=2$.

For $q\ge19$ in this congruence class, fix any arc $u\to v$.
Its common outneighbor set has at least four vertices. Every tournament
on four vertices has a transitive triple: a vertex has at least two
outneighbors, and those two have an arc between them. Such a triple,
together with $u,v$, gives a transitive five-set. Thus $N(q)=3$.

The only prime powers congruent to three modulo four below nineteen are
$3,7,11$; the next is $19$, since $15$ is not a prime power.
The small cases have already been handled. These tournament properties
are standard finite combinatorics; related Paley subtournament calculations
appear in Momihara and Suda~\cite{ms}.

\section{Even characteristic and completion of the proof}
For even $q$, squaring permutes $\F_q^*$, so $\PSL(2,q)=\PGL(2,q)$.
The latter group is sharply triply transitive on $\PP^1(\F_q)$: every
ordered triple of distinct points has a unique projective map to any
other such triple. If $q\ge8$, select three disjoint output pairs.
Every one of their eight choices is a distinct ordered triple and is
therefore realized. Hence $N(q)=3$. The smaller even fields $q=2,4$
were handled in Section~2. This proves Theorem~\ref{thm:main} for every prime power.

\section{Finite verification}
The included standard-library Python verifiers require no network access,
GPU, model API, or paid service. The first constructs finite fields from
irreducible polynomials, enumerates the natural projective action, and
compares cube membership directly with the character criterion.
It checks the fifteen fields
\[
2,3,4,5,7,8,9,11,13,16,17,19,23,25,27.
\]
There are $43\,290$ normalized candidate triples in this sweep. A separate
prime-field implementation uses determinant-one matrices and integer
modular arithmetic for $q=2,3,5,7,11,13,17,19,23$. Counts and hashes of
the full normalized cube sets agree. For example, the counts are $120$
for $q=8$ and $2\,184$ for $q=16$.

These two implementations share the mathematical normalization argument;
they are not independent third-party validations. A user-supplied report
of AI rechecking in another session of the same workflow found no defects
and reported additional calculations without normalization. The associated
custom code and raw logs have not been supplied; that report is not counted
as a reproducible third implementation in this release.
The general proof, rather than extrapolation from this finite sweep,
establishes the all-field claim.

\section{Prior work and limits}
Li~\cite{li} studies learning group elements and Natarajan cubes; several
small cases overlap the present calculation. The known general bounds do
not by themselves give the displayed exception table. No exact all-$q$
classification was located in the checked text and submitted source, but
the reported additional search data were unavailable. This is a limited
literature comparison, not a proof of novelty or priority.

Height, relational complexity, and base size are established invariants
\cite{gls,hudson}; linear dimension of a group action is another, different
quantity~\cite{devillers}. The determinant test, character identities,
Paley structures, and sharp triple transitivity are known ingredients.
For even $q$ the result is an immediate consequence of sharp triple
transitivity and label counting. The point of the note is an explicit
all-field classification for the stated natural action, especially the
odd-characteristic cube criterion and its exceptions, subject to the
unresolved prior-art qualifications.

\noindent\textbf{Disclosure and rights.} AI tools assisted literature checking,
proof drafting, code, and presentation. Checks in this workflow are not
third-party certification. References remain partially checked and novelty
remains inconclusive. Copyright \textcopyright\ 2026 Ryutaro Yonezu.
All rights reserved.

\medskip
\noindent\textbf{Research materials:}\par\nobreak
\noindent{\footnotesize\url{https://github.com/yonezaemon1-hub/no42-psl2-natarajan-dimension}}.

\begin{thebibliography}{9}\small
\bibitem{li} X.~Li.
\emph{Learning a Group Element from Its Action: Galois Correspondences,
Optimal ERM, and an Abel--Ruffini Threshold for the Natarajan Dimension}.
arXiv:2610.18154v1, 2026. \url{https://arxiv.org/abs/2610.18154}.

\bibitem{kpt} V.~Kr\v{c}adinac, M.~O.~Pav\v{c}evi\'{c}, and K.~Tabak.
Three-dimensional Hadamard matrices of Paley type.
\emph{Finite Fields and Their Applications} 92 (2023), 102306.
\url{https://doi.org/10.1016/j.ffa.2023.102306}.

\bibitem{ms} K.~Momihara and S.~Suda.
\emph{Upper bounds on the size of transitive subtournaments in digraphs}.
arXiv:1605.02469, 2016. \url{https://arxiv.org/abs/1605.02469}.

\bibitem{gls} N.~Gill, B.~Lod\`a, and P.~Spiga.
On the height and relational complexity of a finite permutation group.
\emph{Nagoya Mathematical Journal} 246 (2022), 372--411.
\url{https://doi.org/10.1017/nmj.2021.6}.

\bibitem{hudson} S.~Hudson.
\emph{A Paper on Calculating the Height and Relational Complexity of the
Primitive Actions of $\PSL_2(q)$ and $\PGL_2(q)$}.
arXiv:2607.20295v1, 2026. \url{https://arxiv.org/abs/2607.20295}.

\bibitem{devillers} A.~Devillers, M.~Giudici, D.~R.~Hawtin, L.~Klawuhn,
and L.~Morgan. \emph{Linear dimension of group actions}.
arXiv:2512.16079v1, 2025. \url{https://arxiv.org/abs/2512.16079}.
\end{thebibliography}
\end{document}
